\chapter{Blood in Stool}

\section*{Definition}
Blood in your stool means there is bleeding somewhere in your digestive tract. The color of the blood can help tell you where the bleeding is coming from — bright red blood usually means bleeding in the lower part (rectum or colon), while black or tarry stools suggest bleeding higher up (stomach or small intestine).

\section*{When to See a Doctor}
See your doctor if:
\begin{itemize}
  \item Arrange assessment for unexplained or recurrent blood in the stool, particularly with changes in bowel habits, abdominal pain, fatigue, weight loss, or a family history of bowel cancer.
  \item Seek emergency care for non-stop bleeding, a large amount of blood, large clots, faintness, or weakness.
  \item You feel dizzy, lightheaded, or weak (signs of significant blood loss)
  \item Seek urgent care for black, tarry stool, dark-red stool, or bloody diarrhea.
  \item You have abdominal pain, weight loss, or changes in bowel habits along with the bleeding
\end{itemize}

\section*{How It Develops}
Bleeding in the digestive tract can happen for many reasons — from a small tear in the lining (like an anal fissure) to inflammation, infection, or abnormal growths. The appearance of the blood depends on how long it has been in the digestive tract and where it originated. Bright red blood on the toilet paper or coating the stool usually means lower tract bleeding, while dark, tarry stools mean the blood has been digested as it traveled through the intestines.

\section*{Causes}
\begin{itemize}
  \item \textbf{Hemorrhoids:} Swollen veins in the rectum or anus that can bleed, especially during bowel movements
  \item \textbf{Anal fissure:} A small tear in the lining of the anus, often from passing hard stool
  \item \textbf{Diverticulosis:} Small pouches in the colon wall that can bleed
  \item \textbf{Inflammatory bowel disease:} Crohn's disease or ulcerative colitis, which cause inflammation and ulcers in the intestines
  \item \textbf{Peptic ulcer:} A sore in the lining of the stomach or small intestine that can bleed
  \item \textbf{Colorectal polyps or cancer:} Abnormal growths in the colon that may bleed
  \item \textbf{Infections:} Bacterial infections like E. coli or salmonella can cause bloody diarrhea
\end{itemize}

\section*{Related Symptoms}
\begin{itemize}
  \item Abdominal pain or cramping
  \item Change in bowel habits (constipation or diarrhea)
  \item Unexplained weight loss
  \item Fatigue or weakness (from anemia due to blood loss)
  \item Mucus in the stool
\end{itemize}
\textbf{Important:} Blood in the stool has many possible causes, including hemorrhoids and fissures, but it should not be assumed to be benign when it persists, recurs, or occurs with other symptoms.

\section*{Medical Evaluation}
\begin{itemize}
  \item Your doctor will ask about the color and amount of blood, your bowel habits, and any other symptoms you have noticed.
  \item They will do a physical exam, which may include a digital rectal exam to check for hemorrhoids or other issues.
  \item Blood tests may assess anemia and other effects of bleeding. Stool tests may be used for infection or screening when appropriate; a fecal occult-blood test is not a substitute for evaluation of visible blood.
  \item Depending on age, symptoms, examination, and screening history, further assessment may include anoscopy, sigmoidoscopy, colonoscopy, upper endoscopy, or imaging.
\end{itemize}

\section*{Treatment Options}
\begin{itemize}
  \item Treatment depends on the cause. Hemorrhoids or fissures may need fiber, stool-softening measures, topical treatment, or a procedure; antibiotics are used only for selected diagnosed infections.
  \item Inflammatory bowel disease and polyps/cancer require condition-specific medical or specialist treatment.
  \item Polyps found during colonoscopy can usually be removed during the procedure.
  \item Significant bleeding may require emergency procedures to stop the source, including endoscopic treatment or surgery.
\end{itemize}

\section*{Self-Care and Home Management}
\begin{itemize}
  \item Do not rely on hemorrhoid creams to explain or treat unexplained bleeding; ask a clinician or pharmacist before using them.
  \item If constipation or hemorrhoids are contributing, gradually increase dietary fiber and fluids unless a clinician has advised restrictions.
  \item Drink plenty of water and stay hydrated.
  \item Avoid prolonged sitting on the toilet, which can worsen hemorrhoids.
  \item Note: any amount of blood in stool should be discussed with a doctor before relying only on self-care.
\end{itemize}

\section*{References}
\begin{thebibliography}{99}
\bibitem{blood_in_stool:ref1} Jensen DM, Machicado GA. ``Diagnosis and treatment of severe hematochezia.'' \textit{Gastroenterology}. 2018;155(3):623-638.
\bibitem{blood_in_stool:ref2} Gralnek IM, Barkun AN, et al. ``Diagnosis and management of acute lower gastrointestinal bleeding.'' \textit{Lancet}. 2017;389(10068):1923-1934.
\bibitem{blood_in_stool:ref3} Feldman M, Friedman LS, Brandt LJ, eds. \textit{Sleisenger and Fordtran\'s Gastrointestinal and Liver Disease}, 11th ed. Elsevier, 2021.
\bibitem{blood_in_stool:ref4} Zuccaro G. ``Management of the adult patient with acute lower gastrointestinal bleeding.'' \textit{Am J Gastroenterol}. 1998;93(8):1202-1208.
\bibitem{blood_in_stool:ref5} National Health Service. ``Bleeding from the bottom (rectal bleeding).'' Reviewed April 2023. https://www.nhs.uk/symptoms/bleeding-from-the-bottom-rectal-bleeding/.
\end{thebibliography}
